\section{Conclusão}

Os resultados indicam que a granularidade semântica impõe custo mensurável à detecção no TACO, embora não constitua seu único fator limitante. A taxonomia Material representa um compromisso defensável: preserva informação sobre composição e melhora o ponto operacional em relação a Fine, ainda que sua superioridade em AP não seja uniforme. Binary maximiza a detectabilidade ao custo da eliminação dessa informação. A avaliação das mesmas caixas, a decomposição de erros e a análise de cauda longa indicam, em conjunto, a coexistência de confusão semântica, escassez de exemplos e omissões espaciais. Portanto, a taxonomia deve ser selecionada segundo a informação requerida pela aplicação, e não apenas pela maior métrica. O fluxo experimental disponibiliza configurações, mapeamento integral, partição, predições, scripts e artefatos derivados para reprodução e extensão do estudo.
